\chapter*{前言}

这本书的目标不是堆题解，也不是把学习压缩进某个固定周期，而是把数据结构与算法训练成一套可以长期迁移的工程能力。面试题看似零散，背后其实围绕两条主线：一条是\term{数据结构}{data structure}，也就是信息怎样存、怎样查、怎样更新；另一条是\term{算法}{algorithm}，也就是在约束下怎样组织计算、怎样减少重复工作、怎样证明结果正确。数据结构和算法不能分开学。只会说“这题用滑动窗口”，但不知道窗口状态应该放在数组、哈希表还是队列里，代码一定会不稳；只会背 \code{unordered_map}，但不知道它解决的是历史查询问题，也很难在新题里想到它。

本书按“知识完整性”组织，而不是按倒计时组织。前面的章节先建立复杂度、暴力法、复盘方法和 \cpp 代码标准；中间把数组、字符串、链表、栈、队列、哈希表、堆、树、图和并查集讲成数据结构主线；后面再进入双指针、滑动窗口、前缀和、二分、搜索、动态规划、回溯、贪心和综合题。训练周期可以按短期集中训练、学期式推进或长期滚动复习来调整，但正文必须把原理、背景、案例、代码、证明、边界和复盘方法讲完整。

每一类算法都必须从暴力方法开始。暴力方法不是低级答案，而是理解问题空间的方式。两数之和先枚举下标对，才能看出慢点是反复查找补数；最短子数组先枚举所有区间，才能看出正数数组存在窗口单调性；岛屿数量先把每个格子当成节点，才能把题目翻译成图的连通块。优化不是凭感觉，而是来自重复工作、单调性、可维护状态、候选淘汰、最优子结构和数据结构查询能力。

代码使用 \cpp20。示例代码优先使用 \code{const} 引用、明确边界、清楚的变量名和可测试函数。涉及大范围计数、距离、费用、容量、前缀和与几何叉积时，本书统一使用 \code{std::int64_t}；涉及自然溢出的哈希值、位模式和无符号掩码时，统一使用 \code{std::uint64_t}。这些固定宽度类型来自 \code{<cstdint>}，比宽度不透明的整数别名更适合教材讲清边界。书里会解释 \code{vector}、\code{string}、\code{deque}、\code{stack}、\code{queue}、\code{priority_queue}、\code{unordered_map}、\code{set}、树、图和并查集的用法与底层复杂度，也会说明迭代器失效、哈希退化、扩容、拷贝、溢出和递归栈这些真实面试和线上评测里会遇到的问题。
